% Insertar en Conclusiones; no sustituir texto de Cacho.
% 15. Estructura Requerida (Formato LATEX); 19. Distribución de Puntos; 20.3. Criterio 5: Análisis Estadístico (15 pts).
\paragraph{Conclusiones de Othello.}
\textit{Guía: 15. Estructura Requerida (Formato \LaTeX); 20.3. Criterio 5: Análisis Estadístico (15 pts).}
Se implementaron dos evaluaciones acotadas con búsqueda configurable y se verificaron reglas, pases, utilidad terminal, equivalencia con Minimax y registros reproducibles. En las 19 partidas completas, \texttt{fija} ganó 14 y \texttt{etapas} 5, sin empates. La ventaja es descriptiva: los contrastes binomiales por apertura no rechazan la hipótesis nula tras Holm y sólo hay dos o tres bloques por profundidad. Bootstrap y correlaciones se calcularon cuando procedía, con soporte insuficiente para conclusiones firmes; chi-cuadrado no aplica al diseño emparejado.

El costo medido, particularmente a profundidad 7, motivó interrumpir el plan de 104 y cerrar la muestra en 19; las otras 85 no se ejecutaron. No se cumple el criterio de al menos 100 partidas para excelencia estadística. La interrupción deja condiciones desiguales; el subconjunto común de dos aperturas facilita una comparación descriptiva más controlada, pero tampoco demuestra superioridad. Los pesos no están ajustados estadísticamente y la frontera no mide estabilidad exacta. Se incorporaron cinco resultados externos reportados: una victoria y cuatro derrotas con fija a profundidad 4, sin inferencia por condiciones desconocidas y posible dependencia; la terminalidad del último no está confirmada. La participación externa está documentada con evidencia parcial; permanece pendiente integrar el artículo con Cacho. Como trabajo futuro se propone ampliar aperturas bajo un presupuesto computacional previo y realizar ablaciones; esta entrega no afirma haber completado el plan original.
